% !TeX root = ./../../main.tex
% Obwód klasyfikatora z odczytem fazowym dla n = 2 kubitów (K = 3 klasy); rys. fig:cls:circuit (panel a rysunku fig:cls:model), podrozdział 8.2.
% Rysunek własny (quantikz2, zadanie PHD-TS-78), wykonany na podstawie kodu
% /Users/krzysztofwerner/git/quantum-classifier/torus_classifier/circuit.py
% (funkcje _apply_sandwiched_ansatz, ansatz_block, _apply_block_noise).
% Kolejność bramek dwukubitowych w bloku jak w ansatz_block: dla każdej uporządkowanej pary (j, m), j != m,
% bramka sterowana kubitem j, a po niej bramka sterowana kubitem m; dla n = 2 daje to cztery bramki.
% Kanał błędu (linia kropkowana) występuje tylko w symulacji z modelem szumu.
% Każdy kubit jest mierzony w trzech bazach (sigma_x, sigma_y, sigma_z), co daje jego wektor Blocha.
\begin{quantikz}[row sep=0.35cm, column sep=0.13cm, font=\footnotesize]
\lstick{$\ket{0}$} & \gate{\mathrm{H}} & \gate{U}\gategroup[2,steps=6,style={dashed,rounded corners,inner xsep=1pt,inner ysep=3pt},label style={yshift=1pt}]{blok $\ell=1,\dots,L$} & \ctrl{1} & \gate{V} & \gate{V} & \ctrl{1} & \gate[style={densely dotted}]{\errch} & \gate{\mathrm{H}} & \meter{} \\
\lstick{$\ket{0}$} & \gate{\mathrm{H}} & \gate{U} & \gate{V} & \ctrl{-1} & \ctrl{-1} & \gate{V} & \gate[style={densely dotted}]{\errch} & \gate{\mathrm{H}} & \meter{}
\end{quantikz}
